\chapter{Polymorphism and Virtualization}
\label{ch:polymorphism-virtualization}

\textit{One interface, many implementations --- and the machinery that makes it work.}

Inheritance lets you express ``is-a'' relationships and reuse code. Polymorphism elevates that to runtime behaviour: a single base-class pointer can point to any derived object and call the correct implementation automatically. This chapter covers how C++98/03 implements polymorphism through virtual functions and vtables, the hazards that accompany it (the virtual-destructor trap, virtual calls in constructors, the diamond problem), the full taxonomy of inheritance types, \texttt{const} correctness in class hierarchies, alignment and padding in the object model, and the copy-and-swap idiom for exception-safe assignment.

%% ============================================================
\section{Inheritance: The ``Is-A'' Relationship}
\label{sec:inheritance-isa}

\textbf{Inheritance} models an ``is-a'' relationship. A \texttt{Dog} \textit{is an} \texttt{Animal}; a \texttt{Sword} \textit{is a} \texttt{Weapon}. Inheritance lets a \textbf{derived class} acquire the members and methods of a \textbf{base class} without copying code.

\begin{lstlisting}[language=C++, caption={Single public inheritance}]
#include <iostream>

class Animal {
public:
    void eat() { std::cout << "Eating food...\n"; }
};

class Dog : public Animal {
public:
    void bark() { std::cout << "Woof!\n"; }
};

int main() {
    Dog my_dog;
    my_dog.bark();
    my_dog.eat();  // Inherited from Animal
    return 0;
}
\end{lstlisting}

Because a \texttt{Dog} is-a \texttt{Animal}, a \texttt{Dog*} may be passed wherever an \texttt{Animal*} is expected --- this is the \textbf{Liskov Substitution Principle}:

\begin{lstlisting}[language=C++, caption={Passing derived as base}]
void feed_animal(Animal* a) { a->eat(); }

int main() {
    Dog* d = new Dog();
    feed_animal(d); // Valid: Dog is-a Animal
    delete d;
    return 0;
}
\end{lstlisting}

%% ============================================================
\section{The \texttt{protected} Access Modifier}
\label{sec:protected}

\texttt{protected} members are \textbf{private to the outside world} but \textbf{fully accessible to derived classes}:

\begin{lstlisting}[language=C++, caption={protected member shared with derived classes}]
class Animal {
protected:
    int weight;
};

class Dog : public Animal {
public:
    void grow() { weight += 5; }
};
\end{lstlisting}

%% ============================================================
\section{Inheritance Specifiers: \texttt{public}, \texttt{protected}, \texttt{private}}
\label{sec:inheritance-specifiers}

\begin{center}
\begin{tabular}{|l|l|l|l|}
\hline
\textbf{Base member} & \textbf{public inh.} & \textbf{protected inh.} & \textbf{private inh.} \\
\hline
\texttt{public}    & \texttt{public}    & \texttt{protected} & \texttt{private} \\
\texttt{protected} & \texttt{protected} & \texttt{protected} & \texttt{private} \\
\texttt{private}   & inaccessible       & inaccessible       & inaccessible     \\
\hline
\end{tabular}
\end{center}

\begin{lstlisting}[language=C++, caption={Access under three inheritance modes}]
struct A {
public:    int p1;
protected: int p2;
private:   int p3;
};

struct B_pub : public A {
    void foo() { p1 = 0; p2 = 0; /* p3 inaccessible */ }
};

struct B_pri : private A {
    void foo() { p1 = 0; p2 = 0; /* both private in B_pri */ }
};
\end{lstlisting}

\textbf{Public inheritance} models ``is-a''. \textbf{Private inheritance} models ``is-implemented-in-terms-of''. \textbf{Protected inheritance} is rarely used.

%% ============================================================
\section{Virtual Functions and Dynamic Dispatch}
\label{sec:virtual-functions}

Without virtual functions, method calls on a base-class pointer are resolved at \textbf{compile time} (static dispatch):

\begin{lstlisting}[language=C++, caption={Non-virtual --- always calls base version}]
#include <iostream>
struct X { void f() { std::cout << "X::f()\n"; } };
struct Y : X { void f() { std::cout << "Y::f()\n"; } };

void call(X& a) { a.f(); }

int main() {
    X x; Y y;
    call(x); // X::f()
    call(y); // X::f() --- Y::f() is NEVER called!
    return 0;
}
\end{lstlisting}

Marking the base function \texttt{virtual} switches to \textbf{dynamic dispatch}:

\begin{lstlisting}[language=C++, caption={virtual --- correct override called at runtime}]
#include <iostream>

class Animal {
public:
    virtual void speak() { std::cout << "...\n"; }
    virtual ~Animal() {}
};

class Dog : public Animal {
public:
    void speak() { std::cout << "Woof!\n"; }
};

class Cat : public Animal {
public:
    void speak() { std::cout << "Meow!\n"; }
};

int main() {
    Animal* pet = new Dog();
    pet->speak(); // "Woof!" -- resolved at runtime
    delete pet;
    return 0;
}
\end{lstlisting}

%% ============================================================
\section{The vTable Mechanism}
\label{sec:vtable}

When a class contains at least one virtual function, the compiler:

\begin{enumerate}
    \item Creates a \textbf{vtable} (virtual table) for the class --- a static array of function pointers, one entry per virtual function.
    \item Inserts a hidden \textbf{vptr} (virtual pointer) into every instance, pointing to its class's vtable.
\end{enumerate}

\begin{lstlisting}[language=C++, caption={vtable overhead illustration}]
class Base {
    int data;
    virtual void func() {}
};
// On 64-bit: sizeof(Base) = sizeof(vptr=8) + sizeof(int=4) + padding = 16 bytes
// Without virtual: sizeof(Base) = sizeof(int=4) = 4 bytes
\end{lstlisting}

Dynamic dispatch adds one pointer dereference per virtual call. In tight loops this can cause cache misses; this is a real concern in game engines and HFT hot paths.

%% ============================================================
\section{Pure Virtual Functions and Abstract Classes}
\label{sec:pure-virtual}

A \textbf{pure virtual function} (\texttt{= 0}) has no default implementation. Any class with at least one pure virtual function is \textbf{abstract} --- it cannot be instantiated.

\begin{lstlisting}[language=C++, caption={Abstract base class with pure virtual functions}]
#include <iostream>

class Weapon {
public:
    virtual void attack() = 0;
    virtual ~Weapon() {}
};

class Sword : public Weapon {
public:
    void attack() { std::cout << "Swing!\n"; }
};

int main() {
    // Weapon w; // ERROR: cannot instantiate abstract class
    Weapon* w = new Sword();
    w->attack();
    delete w;
    return 0;
}
\end{lstlisting}

A pure virtual function may still have a body --- derived classes can call it via \texttt{Base::func()}:

\begin{lstlisting}[language=C++, caption={Pure virtual with default body}]
struct DefaultAbstract {
    virtual void configure() = 0;
};

void DefaultAbstract::configure() { /* shared setup */ }

struct Derived : DefaultAbstract {
    void configure() { DefaultAbstract::configure(); }
};
\end{lstlisting}

%% ============================================================
\section{The Virtual Destructor Trap}
\label{sec:virtual-destructor}

\begin{lstlisting}[language=C++, caption={Non-virtual destructor causes resource leak}]
class Base {
public:
    ~Base() { std::cout << "Base destroyed.\n"; }
};

class Derived : public Base {
    int* array;
public:
    Derived() { array = new int[100]; }
    ~Derived() { delete[] array; std::cout << "Derived destroyed.\n"; }
};

int main() {
    Base* b = new Derived();
    delete b; // ONLY prints "Base destroyed." -- array is LEAKED!
    return 0;
}
\end{lstlisting}

\textbf{The Golden Rule:} if a class has even one \texttt{virtual} function, give it a \texttt{virtual} destructor:

\begin{lstlisting}[language=C++, caption={Correct --- virtual destructor}]
class Base {
public:
    virtual ~Base() { std::cout << "Base destroyed.\n"; }
};
\end{lstlisting}

To prevent deletion through a base pointer entirely, make the destructor \texttt{protected} (non-virtual):

\begin{lstlisting}[language=C++, caption={Protected destructor blocks base-pointer deletion}]
struct NoDeletionViaBase {
protected:
    ~NoDeletionViaBase() {}
};
\end{lstlisting}

%% ============================================================
\section{Virtual Functions in Constructors and Destructors}
\label{sec:virtual-in-ctor}

Never call virtual functions during construction or destruction:

\begin{lstlisting}[language=C++, caption={Virtual call during construction does not dispatch to derived}]
#include <iostream>
using namespace std;

class Base {
public:
    Base()  { f("base constructor"); }
    ~Base() { f("base destructor"); }
    virtual const char* v() { return "Base::v()"; }
    void f(const char* caller) {
        cout << "From " << caller << ": " << v() << "\n";
    }
};

class Derived : public Base {
public:
    Derived()  { f("derived constructor"); }
    ~Derived() { f("derived destructor"); }
    const char* v() { return "Derived::v()"; }
};

int main() {
    Derived d;
    return 0;
}
/* Output:
   From base constructor: Base::v()    <-- NOT Derived::v()!
   From derived constructor: Derived::v()
   From derived destructor: Derived::v()
   From base destructor: Base::v()     <-- NOT Derived::v()! */
\end{lstlisting}

During \texttt{Base::Base()}, the dynamic type is \texttt{Base}; the derived subobject does not yet exist.

%% ============================================================
\section{Multiple Inheritance and Thunks}
\label{sec:multiple-inheritance}

\begin{lstlisting}[language=C++, caption={Multiple inheritance and subobject layout}]
class A { int a; virtual void f() {} };
class B { int b; virtual void g() {} };
class C : public A, public B { int c; };

int main() {
    C obj;
    A* pa = &obj; // Points to start of obj (A subobject at offset 0)
    B* pb = &obj; // Points to obj + sizeof(A): different address!
    return 0;
}
\end{lstlisting}

When a virtual function from \texttt{B} is called through \texttt{pb}, the compiler generates a \textbf{thunk} --- a small code fragment --- to adjust the \texttt{this} pointer before entering the function body.

%% ============================================================
\section{Virtual Inheritance and the Diamond Problem}
\label{sec:diamond-problem}

Without virtual inheritance, two base classes that share a common ancestor duplicate it:

\begin{lstlisting}[language=C++, caption={Diamond problem --- ambiguous Top::t}]
struct Top { int t; };
struct Left  : public Top { int l; };
struct Right : public Top { int r; };
struct Bottom : public Left, public Right { int b; };

void f() {
    Bottom obj;
    // obj.t;         // ERROR: ambiguous
    obj.Left::t = 1;  // Two separate copies of t
    obj.Right::t = 2;
}
\end{lstlisting}

\texttt{virtual} inheritance ensures one shared instance:

\begin{lstlisting}[language=C++, caption={Virtual inheritance resolves the diamond}]
struct Top  { int t; };
struct Left  : virtual public Top { int l; };
struct Right : virtual public Top { int r; };
struct Bottom : public Left, public Right { int b; };

void f() {
    Bottom obj;
    obj.t = 42; // Unambiguous: one copy of Top
}
\end{lstlisting}

\textbf{Cost:} \texttt{Left} and \texttt{Right} each contain a hidden \textbf{vbptr} (virtual base pointer). Accessing \texttt{Top} members requires an extra indirection. The most-derived class (\texttt{Bottom}) must directly initialise \texttt{Top}.

%% ============================================================
\section{Multilevel and Hierarchical Inheritance}
\label{sec:inheritance-types}

\subsection{Multilevel Inheritance}

\begin{lstlisting}[language=C++, caption={Three-level inheritance chain}]
#include <iostream>
#include <string>
using namespace std;

class Vehicle {
protected:
    string brand;
public:
    Vehicle(const string& b) : brand(b) {}
    virtual void start() { cout << brand << " is starting\n"; }
    virtual ~Vehicle() {}
};

class Car : public Vehicle {
protected:
    int numDoors;
public:
    Car(const string& b, int doors) : Vehicle(b), numDoors(doors) {}
    void start() {
        cout << brand << " car (" << numDoors << " doors) starting\n";
    }
    virtual ~Car() {}
};

class ElectricCar : public Car {
    int batteryPct;
public:
    ElectricCar(const string& b, int doors, int bat)
        : Car(b, doors), batteryPct(bat) {}
    void start() {
        cout << "Electric " << brand << " (battery: "
             << batteryPct << "%) starting\n";
    }
};

int main() {
    ElectricCar tesla("Tesla", 4, 95);
    tesla.start();
    return 0;
}
\end{lstlisting}

\subsection{Hierarchical Inheritance}

\begin{lstlisting}[language=C++, caption={Hierarchical inheritance --- three employee types}]
#include <iostream>
#include <string>
using namespace std;

class Employee {
protected:
    string name;
    double salary;
public:
    Employee(const string& n, double s) : name(n), salary(s) {}
    virtual void work() = 0;
    virtual ~Employee() {}
};

class Manager : public Employee {
public:
    Manager(const string& n, double s) : Employee(n, s) {}
    void work() { cout << name << " is managing the team\n"; }
};

class Developer : public Employee {
    string language;
public:
    Developer(const string& n, double s, const string& lang)
        : Employee(n, s), language(lang) {}
    void work() { cout << name << " is coding in " << language << "\n"; }
};

int main() {
    Manager   alice("Alice", 80000);
    Developer bob("Bob", 70000, "C++");
    alice.work();
    bob.work();
    return 0;
}
\end{lstlisting}

%% ============================================================
\section{Polymorphic Design: The Shape Pattern}
\label{sec:shape-pattern}

\begin{lstlisting}[language=C++, caption={Polymorphic Shape hierarchy}]
#include <iostream>
#include <string>
using namespace std;

class Shape {
public:
    virtual void   draw()   = 0;
    virtual double area()   = 0;
    virtual string name()   = 0;
    virtual ~Shape() {}
};

class Circle : public Shape {
    double radius;
public:
    explicit Circle(double r) : radius(r) {}
    void   draw()   { cout << "Drawing Circle\n"; }
    double area()   { return 3.14159 * radius * radius; }
    string name()   { return "Circle"; }
};

class Rectangle : public Shape {
    double w, h;
public:
    Rectangle(double w, double h) : w(w), h(h) {}
    void   draw()   { cout << "Drawing Rectangle\n"; }
    double area()   { return w * h; }
    string name()   { return "Rectangle"; }
};

int main() {
    Shape* shapes[2];
    shapes[0] = new Circle(5.0);
    shapes[1] = new Rectangle(4.0, 6.0);

    for (int i = 0; i < 2; ++i) {
        shapes[i]->draw();
        cout << "Area: " << shapes[i]->area() << "\n";
        delete shapes[i];
    }
    return 0;
}
\end{lstlisting}

%% ============================================================
\section{Safe Downcasting in Polymorphic Hierarchies}
\label{sec:safe-downcast}

\begin{lstlisting}[language=C++, caption={Safe downcast with dynamic\_cast}]
#include <iostream>
using namespace std;

class Shape {
public:
    virtual double area() const = 0;
    virtual ~Shape() {}
};

class Circle : public Shape {
    double radius;
public:
    explicit Circle(double r) : radius(r) {}
    double area()     const { return 3.14159 * radius * radius; }
    double diameter() const { return 2.0 * radius; }
};

int main() {
    Shape* s = new Circle(5.0);
    cout << "Area: " << s->area() << "\n";

    Circle* c = dynamic_cast<Circle*>(s);
    if (c != NULL) {
        cout << "Diameter: " << c->diameter() << "\n";
    }

    delete s;
    return 0;
}
\end{lstlisting}

%% ============================================================
\section{\texttt{const} Correctness in OOP}
\label{sec:const-correctness}

\textbf{const member functions} promise not to modify any non-\texttt{mutable} member. They are callable on \texttt{const} instances; non-\texttt{const} methods are not.

\begin{lstlisting}[language=C++, caption={const member functions and mutable}]
#include <iostream>
#include <string>
using namespace std;

class Person {
    mutable int accessCount;
    string name;
    int age;
public:
    Person(const string& n, int a) : name(n), age(a), accessCount(0) {}

    string getName() const {
        ++accessCount; // OK: mutable
        return name;
    }

    void setAge(int a) { age = a; } // Non-const

    int getAge() const { return age; }
};

int main() {
    const Person cp("Bob", 25);
    cout << cp.getName() << "\n";
    cout << cp.getAge()  << "\n";
    // cp.setAge(26); // ERROR: non-const method on const object
    return 0;
}
\end{lstlisting}

%% ============================================================
\section{Alignment, Padding, and the Cost of Polymorphism}
\label{sec:alignment-padding}

A member of size $N$ must reside at an offset divisible by $N$. The compiler inserts \textbf{padding} bytes to satisfy this constraint.

\begin{lstlisting}[language=C++, caption={Struct padding}]
struct Mixed {
    char  a;   // 1 byte at offset 0
               // 3 bytes padding
    int   b;   // 4 bytes at offset 4
    short c;   // 2 bytes at offset 8
               // 6 bytes padding
};
// sizeof(Mixed) = 16 on 64-bit
\end{lstlisting}

Sort members largest-to-smallest to minimise padding:

\begin{lstlisting}[language=C++, caption={Padding-minimised layout}]
struct Optimised {
    int   b;   // 4 bytes at offset 0
    short c;   // 2 bytes at offset 4
    char  a;   // 1 byte at offset 6
               // 1 byte padding
};
// sizeof(Optimised) = 8
\end{lstlisting}

Every polymorphic class adds a hidden vptr (8 bytes on 64-bit) plus alignment padding:

\begin{lstlisting}[language=C++, caption={vptr adds size to the object}]
class WithVirtual {
    int data;
    virtual void f() {}
};
// sizeof(WithVirtual) = 16 (vptr=8, int=4, padding=4) on 64-bit
\end{lstlisting}

%% ============================================================
\section{The Copy-and-Swap Idiom}
\label{sec:copy-and-swap}

The naive copy-assignment operator violates the strong exception guarantee and duplicates copy-constructor logic. \textbf{Copy-and-swap} solves both:

\begin{lstlisting}[language=C++, caption={Copy-and-swap for exception-safe assignment}]
#include <cstring>
#include <algorithm>

class Person {
    char* name;
    int   age;
public:
    Person(const char* n, int a)
        : name(new char[std::strlen(n) + 1]), age(a) {
        std::strcpy(name, n);
    }

    ~Person() { delete[] name; }

    Person(const Person& other)
        : name(new char[std::strlen(other.name) + 1]), age(other.age) {
        std::strcpy(name, other.name);
    }

    friend void swap(Person& lhs, Person& rhs) {
        std::swap(lhs.name, rhs.name);
        std::swap(lhs.age,  rhs.age);
    }

    Person& operator=(Person rhs) { // rhs is a copy
        swap(*this, rhs);
        return *this;
    } // rhs (old state) is destroyed here
};
\end{lstlisting}

If copy-construction throws, the assignment operator body is never entered and \texttt{*this} is untouched --- the strong exception guarantee is satisfied.

%% ============================================================
\section{Professional Insights: Virtual Mechanics and Destructor Policy}
\label{sec:pro-insights-ch6}

\subsection{Virtual Destructor Policy}

\begin{center}
\begin{tabular}{|l|l|}
\hline
\textbf{Scenario} & \textbf{Required destructor} \\
\hline
Polymorphic base class                & \texttt{virtual \~{}Base()} \\
Prevent base-pointer deletion         & \texttt{protected \~{}Base()} \\
Concrete leaf class, no polymorphism  & Non-virtual (implicit or explicit) \\
Pure interface (all pure virtual)     & \texttt{virtual \~{}Interface() = 0} + out-of-line body \\
\hline
\end{tabular}
\end{center}

\subsection{\texttt{override} and \texttt{final} --- Forward Reference: C++11}

C++11 introduced two contextual keywords:

\begin{itemize}
    \item \textbf{\texttt{override}} --- instructs the compiler to verify the function actually overrides a base virtual function. A signature mismatch becomes a compile error instead of a silent new overload.
    \item \textbf{\texttt{final}} --- prevents further overriding or further subclassing.
\end{itemize}

In C++98/03, a mismatched signature silently creates a new overload --- a notoriously hard-to-find bug.

\subsection{Never Call Virtual Functions in Constructors or Destructors}

During a base-class constructor, the dynamic type is the base class. Calling a virtual function calls the base implementation, not the derived override. Design constructors to avoid virtual dispatch.

\subsection{The Liskov Substitution Principle}

Public inheritance is correctly used only when an instance of the derived class can be substituted for the base in every context without breaking correctness. Private inheritance (``is-implemented-in-terms-of'') should be used when you want the implementation but not the interface of the base class.
